\chapter{Literature Review}
\label{ch:literature_review}

This chapter covers the existing research on machine learning for disease prediction. I start with the general use of ML in healthcare, then narrow down to symptom-based prediction, discuss the main challenges, and end with the healthcare booking context that motivates this work.

\section{Machine Learning in Healthcare}
\label{sec:ml_healthcare}

Machine learning has been applied to healthcare for over a decade, but the field has accelerated with the digitization of medical records and cheaper computing power. \citet{Topol2019} provided an early overview of how AI is changing clinical practice, showing that ML models can match or exceed human performance in specific tasks like reading X-rays or pathology slides. \citet{Rajpurkar2022} later surveyed the broader landscape and pointed out that while image-based diagnosis has received most of the attention, structured clinical data --- such as symptom records --- remains underexplored.

In healthcare, the dominant ML approach is supervised classification: train a model on labeled examples (patient features + diagnosis) and use it to predict diagnoses for new patients. The classifiers used most often include Naive Bayes, Logistic Regression, SVM, Decision Trees, Random Forest, and XGBoost \citep{Uddin2019}. \citet{Sarker2021} provides a useful categorization of these methods.

\citet{Davenport2019} divided healthcare AI into three areas: clinical decision support, patient engagement, and administrative efficiency. Disease prediction falls into the first category. The key point is that these models are meant to assist doctors, not replace them. This means they need to be not just accurate but also interpretable and safe.

One thing that stood out to me while reading the literature is that accuracy alone is not a good enough metric for medical ML. \citet{Rajpurkar2022} argued that precision, recall, and AUC-ROC give a better picture, especially on imbalanced datasets where rare diseases can be completely missed by an accuracy-optimized model. This is why I adopted a multi-metric approach in my evaluation (Chapter~\ref{ch:results}).

\section{Symptom-Based Disease Prediction}
\label{sec:symptom_prediction}

Symptom-based disease prediction is a specific type of supervised classification where the inputs are patient-reported symptoms and the output is a disease category. Compared to image-based or genomics-based approaches, symptom-based prediction uses simpler features --- typically binary indicators for whether each symptom is present or absent. This simplicity makes it practical for patient-facing applications and low-resource settings.

Several studies have tested individual classifiers on symptom data. \citet{Uddin2019} compared Naive Bayes, Decision Trees, and Random Forest on multiple medical datasets and found that ensemble methods generally performed better, though the datasets and preprocessing varied between experiments, making it hard to draw firm conclusions.

\citet{Chowdhury2024} reviewed the ML-for-diagnosis literature systematically and found that while Random Forest and XGBoost consistently ranked high, very few studies compared multiple classifiers under the same conditions. They specifically called for controlled comparisons with identical preprocessing and evaluation --- which is exactly what I do in this thesis.

Most studies represent symptoms as binary vectors, where each position corresponds to a symptom and the value is 1 (present) or 0 (absent). This is simple and effective, though it ignores severity and duration. Despite this limitation, binary vectors work well for disease classification, especially with ensemble methods that can pick up on feature interactions \citep{Breiman2001}.

One thing I noticed while reviewing the literature is that almost all studies use a single dataset. This makes it impossible to tell whether the results generalize or are specific to that particular data. I decided to test on two datasets of very different sizes to address this: one with about 5,000 records \citep{Kaggle_Disease2020} and one with about 246,000 \citep{Kaggle_Diseases2023}.

\section{Challenges in Medical Prediction}
\label{sec:challenges}

Medical classification has challenges that standard ML benchmarks do not. The biggest one is class imbalance: some diseases are common and others are rare, so the training data is skewed. A classifier can achieve high accuracy by simply predicting the most common diseases and ignoring the rare ones \citep{Haixiang2017}.

\citet{Chawla2002} proposed SMOTE (Synthetic Minority Over-sampling Technique), which creates synthetic examples of rare classes by interpolating between existing ones. SMOTE has been widely used, but it has problems with binary data --- interpolating between two binary vectors can produce unrealistic symptom combinations \citep{Fernandez2018}. An alternative is class weighting, where the loss function penalizes misclassifications of rare diseases more heavily. I chose class weighting for this thesis because it works naturally with binary features and doesn't create artificial symptom patterns.

\citet{Thabtah2020} tested multiple imbalance-handling techniques across different classifiers and found that no single technique works best everywhere --- it depends on the classifier and the degree of imbalance. This is one of the reasons I test five different classifiers rather than just picking one.

Beyond imbalance, symptom data is noisy. Patients may forget to mention symptoms, describe them imprecisely, or report symptoms that are not relevant to their condition \citep{Chowdhury2024}. Public datasets, while useful for research, may not fully capture the messiness of real clinical encounters.

\section{Healthcare Booking Context}
\label{sec:booking_context}

The practical motivation behind this work is healthcare booking. When patients feel sick, they have to decide whom to see: a general practitioner, a specialist, or an emergency department. Most people make this decision with limited medical knowledge, which leads to wasted appointments, unnecessary referrals, and delayed treatment for serious conditions.

A symptom-based prediction model could help here. If a patient enters their symptoms into a booking system and the model predicts a dermatological condition, the system could suggest booking directly with a dermatologist. The WHO's digital health strategy supports this kind of tool for improving healthcare access \citep{WHO2021}.

But deploying ML in healthcare raises serious concerns. The EU's General Data Protection Regulation \citep{GDPR2016} has strict rules about processing health data, and the new EU AI Act \citep{EUAI2024} classifies medical AI systems as high-risk, requiring documented risk assessments and human oversight. I discuss these regulatory aspects in more detail in Chapter~\ref{ch:discussion}.
